\documentclass[10pt,a4paper]{amsart}
\usepackage[margin=19mm]{geometry}
\usepackage{amsmath,amssymb,mathtools}
\usepackage[scr=rsfs,cal=euler]{mathalfa}
\usepackage{xfrac}
\usepackage[unicode,hidelinks,bookmarks=false]{hyperref}
\newcommand{\ie}{i.e.\ }
\newcommand{\cf}{cf.\ }
\newcommand{\resp}{respectively\ }
\newcommand{\Subsection}{\subsection}
\providecommand{\nobreakdash}{\nobreak\hbox{-}}
\setlength{\emergencystretch}{2em}
\multlinegap0pt
\usepackage{bm}
\usepackage{bbm}
\usepackage{dsfont}
\usepackage{xspace}
\usepackage{cancel}
\usepackage{mathtools}
\usepackage{amsthm}
\makeatletter
\@ifundefined{thm}{\newtheorem{thm}{Thm}}{}
\@ifundefined{defn}{\newtheorem{defn}[thm]{Definition}}{}
\@ifundefined{prop}{\newtheorem{prop}[thm]{Proposition}}{}
\makeatother
\title[Satellite knots and immersed Heegaard Floer homology]{Satellite knots and immersed Heegaard Floer homology}
\author{Wenzhao Chen, Jonathan Hanselman}
\date{Source: arXiv:2309.12297v1 (CC BY 4.0)}
\begin{document}
\maketitle

\noindent\emph{Source and licence.} Satellite knots and immersed Heegaard Floer homology. Version arXiv:2309.12297v1 (CC BY 4.0), distributed under the Creative Commons Attribution 4.0 International licence, as recorded in the arXiv licence metadata for this exact version and shown on the abstract page https://arxiv.org/abs/2309.12297v1. Full rights notice: licenses/arxiv-2309.12297-CC-BY-4.0.txt.\par\smallskip

\noindent\textit{Editorial scope.} Counted unit: the Prop whose source label is \texttt{Proposition, weakly extended type D strucuture is well defined.} in the article (source0.tex lines 1085--1087, source label \texttt{Proposition, weakly extended type D strucuture is well defined.}), delivered together with the article's own complete original proof of it (source0.tex lines 1088--1133, one \texttt{proof} environment of 46 lines). No mathematics is written, repaired or re-derived: statement and proof are the article's own text line for line. Required context retained uncounted: Proposition, ends of 1-dimensional moduli spaces of 0-P curves, Reeb element of length 4 (lines 707--722); Proposition, ends of moduli space of 1-P curves (lines 899--913); Definition, type D structure (lines 1035--1039); Proposition, type D structure is well defined (lines 1046--1048); Definition, weakly extended type D structure (lines 1075--1079). Every definition, display and label of those lines is retained unchanged. Omitted from this package and not redistributed: 1--706, 723--898, 914--1034, 1040--1045, 1049--1074, 1080--1084, 1134--2332. Nothing inside a retained excerpt is omitted or altered: every retained body is delivered exactly as the article's lines are written, so no omission receipt of any kind accompanies this package. Citations: the retained excerpts carry no citation command at all, so no bibliography entry of the article is transcribed and no reference is redistributed. Reference adaptation: none was needed and none was made; every reference inside the delivered unit resolves inside it, so no unresolved reference or citation is produced. Printed slips kept literal: no printed word, symbol, hypothesis or number of the counted statement or of its proof was altered. Layout and class: the article's own class is \texttt{amsart} with packages including enumerate, tikz, bm, hyperref, mathrsfs, amsfonts, latexsym, epsfig, indentfirst, amsmath, amsthm, graphicx, amssymb, amscd; the wrapper is the lane's pinned \texttt{amsart} editorial wrapper at 10pt on A4, which restates the theorem-like environments the retained text needs and adds the article's own macro definitions that the retained text needs (none), with extra wrapper packages (bm, bbm, dsfont, xspace, cancel, mathtools). The delivered numbering is the unit's own. Assets: no retained span contains a figure environment, a graphics-inclusion command, a PDF-page inclusion, a table or a TikZ picture; no image file is copied into this package, the package declares no asset dependency and no image file is redistributed with it. Licence: version arXiv:2309.12297v1 of this article is distributed under the Creative Commons Attribution 4.0 International licence, as recorded in the arXiv licence metadata for this exact version and shown on the abstract page https://arxiv.org/abs/2309.12297v1. A full rights notice is delivered at licenses/arxiv-2309.12297-CC-BY-4.0.txt. Characters: every retained excerpt is pure ASCII, so the delivered engine, XeLaTeX with its default Latin Modern text font, reports no missing character.
\begin{prop}\label{Proposition, ends of 1-dimensional moduli spaces of 0-P curves, Reeb element of length 4}
Given a compatible pair $(B,\overrightarrow{\bold{\rho}})$ with $B\in \pi_2(\bm{x},\bm{y})$, $\overrightarrow{\rho}\in\{(-\rho_0,-\rho_1,-\rho_2,-\rho_3),(-\rho_0,-\rho_{123}),(-\rho_{012},-\rho_3)\}$, and $\text{Ind}(B,\overrightarrow{\bold{\rho}})=2$, the compactified moduli space $\mathcal{M}^B(\bm{x},\bm{y};\overrightarrow{\bold{\rho}})$ is a compact 1-manifold with boundary. The ends are of the following four types:
\begin{itemize}
\item[(E-1)]Two-story ends;
\item[(E-2)]split curve ends;
\item[(E-3)]ends corresponding to boundary degeneration with corners;
\item[(E-4)]ends corresponding to boundary degeneration without corners.
\end{itemize}
Moreover, 
\begin{itemize}
\item[(a)]If split curve ends occur, then $\overrightarrow{\bold{\rho}}$ is either $(-\rho_0,-\rho_{123})$ or $(-\rho_{012},-\rho_3)$, and the number of such ends is $\#\mathcal{M}^B(\bm{x},\bm{y};-\rho_{1230})$ if $\overrightarrow{\bold{\rho}}=(-\rho_0,-\rho_{123})$; otherwise the number of ends is equal to $\#\mathcal{M}^B(\bm{x},\bm{y};-\rho_{3012})$;
\item[(b)]If ends corresponding to boundary degeneration with corners occur and $\overrightarrow{\rho}=(-\rho_0,-\rho_1,-\rho_2,-\rho_3)$, then the number of such ends is mod 2 congruent to $$\sum_{\{(B_1,\ q)|\exists B_2\in T(q),\ B_1+B_2=B\}}\#\mathcal{M}^{B_1}(\bm{x},\bm{y};q).$$
For the other two choices of $\overrightarrow{\rho}$, the numbers of ends corresponding to boundary degeneration with corners are both even;
\item[(c)]If (E-4) occurs, then $[B]=[\Sigma]$, $\bm{x}=\bm{y}$, and the number of such ends is odd.
\end{itemize}
\end{prop}

\begin{prop}\label{Proposition, ends of moduli space of 1-P curves}
Let $B\in\tilde{\pi}_2(\bm{x},\bm{y})$ such that $\iota(\bm{x})=\iota_1$ and $\text{ind}(B;U)=2$. Then fixing a generic almost complex structure, the compactified moduli space $\overline{\mathcal{M}}^B(\bm{x},\bm{y};U)$ is a compact 1-manifold with boundary. The boundaries are of the following types:
\begin{itemize}
\item[(1)]Two-story building
\item[(2)]simple holomorphic combs $(u,v)$ with $v$ being an orbit curve
\item[(3)]boundary degeneration with corners
\item[(4)]boundary degeneration without corners
\end{itemize}
Moreover, 
\begin{itemize}
\item[(a)]The number of type (2) ends is $\#\mathcal{M}^B(\bm{x},\bm{y};-\rho_{1230})+\#\mathcal{M}^B(\bm{x},\bm{y};-\rho_{3012})$.
\item[(b)]The number of type (3) ends is mod 2 congruent to $$\sum_{\{(B_1,\ q)|B_2\in T(q),\ B_1+B_2=B\}}\#\mathcal{M}^{B_1}(\bm{x},\bm{y};q)$$
\item[(c)]The number of type (4) ends is even.
\end{itemize}
\end{prop}

\begin{defn}\label{Definition, type D structure}
Let $\mathcal{H}$ be an unobstructed, provincially admissible, immersed bordered Heegaard diagram. Fix a generic almost complex structure on $\Sigma\times [0,1] \times \mathbb{R}$. The type D module $\widehat{CFD}(\mathcal{H})$ is defined to be the $\mathcal{A}$-module $$\mathcal{A}\otimes_{\mathcal{I}}X^\mathcal{E}(\mathcal{H})$$ together with a differential given by
$$\partial (a\otimes \eta)=a\cdot(\sum_{\bm{y}}\sum_{\{(B,\overrightarrow{\sigma})|\ n_z(B)=0,\ \text{ind}(B,\overrightarrow{\sigma})=1\}}\#\mathcal{M}^B(\bm{x},\bm{y};\overrightarrow{\sigma})a(-\overrightarrow{\sigma})\otimes \Phi^B_{\bm{x},\bm{y}}\eta),$$
where $a\in\mathcal{A}$, $\eta\in \mathcal{E}|_{\bm{x}}$, and the pairs $(B,\overrightarrow{\sigma})$ are compatible. The underlying type D structure is the pair $(X^\mathcal{E}(\mathcal{H}), \delta)$ where $\delta(\eta)\coloneqq \partial(1\otimes \eta)$ for any $\eta\in X^\mathcal{E}(\mathcal{H})$.
\end{defn}

\begin{prop}\label{Proposition, type D structure is well defined}
The operator $\partial$ in Definition \ref{Definition, type D structure} is well-defined and $\partial^2=0$. 
\end{prop}

\begin{defn}\label{Definition, weakly extended type D structure}
Let $\mathcal{H}$ be an unobstructed, provincially admissible, immersed bordered Heegaard diagram. Fix a generic admissible almost complex structure on $\Sigma\times [0,1] \times \mathbb{R}$. The weakly extended type D module $\widetilde{CFD}(\mathcal{H})$ is defined to be the $\tilde{\mathcal{A}}$-module $$\tilde{\mathcal{A}}\otimes_{\mathcal{I}}X^{\mathcal{E}}(\mathcal{H})$$ together with a differential given by 
$$\tilde{\partial} (a\otimes \eta)=a\cdot(\sum_{\bm{y}}\sum_{\{(B,\overrightarrow{\sigma})|\ \text{ind}(B,\overrightarrow{\sigma})=1\}}\#\mathcal{M}^B(\bm{x},\bm{y};\overrightarrow{\sigma})a(-\overrightarrow{\sigma})\otimes \Phi^B_{\bm{x},\bm{y}}\eta),$$
where $a\in\tilde{\mathcal{A}}$, $\eta\in \mathcal{E}|_{\bm{x}}$, $\overrightarrow{\sigma}$ is a sequence of Reeb chords that include the case of a single closed Reeb orbit $\{U\}$ (in which case the corresponding moduli space consists of 1-P holomorphic curves), and the pairs $(B,\overrightarrow{\sigma})$ are compatible. When $\overrightarrow{\sigma}=\{U\}$, we define $a(-U)=\bm{U}$. The underlying weakly extended type D structure is $(X^{\mathcal{E}}(\mathcal{H}),\tilde{\delta})$ where $\tilde{\delta}(\eta)\coloneqq \tilde{\partial}(1\otimes \eta)$.
\end{defn}

\begin{prop}\label{Proposition, weakly extended type D strucuture is well defined.}
The operator $\tilde{\partial}$ in Definition \ref{Definition, weakly extended type D structure} is well-defined and $\tilde{\partial}^2=\bm{U}$. 
\end{prop}

\begin{proof}
A standard argument shows that the provincial admissibility of $\mathcal{H}$ implies the sum defining $\tilde{\partial}$ in Definition \ref{Definition, weakly extended type D structure} is finite, and hence $\tilde{\partial}$ is well-defined. 

Next, we show $\tilde{\partial}^2=\bm{U}$. Once again, we first give the proof when the local system is trivial for conciseness. Recall the length of an element $a\in\tilde{\mathcal{A}}$ is the number of factors $\rho_i\in\{\rho_0,\rho_1,\rho_2,\rho_3\}$ when we write $a$ as a product of the generators $\{\iota_0,\iota_1,\rho_0,\rho_1,\rho_2,\rho_3\}$. (For example, $\rho_{123}$ has length $3$ and $\iota_0$ has length $0$.) 

For an element $a\in\tilde{\mathcal{A}}$ whose length is less than or equal to $3$, the proof of Proposition \ref{Proposition, type D structure is well defined} carries over to show $\langle \tilde{\partial}^2\bm{x}, a\bm{y}\rangle =0$ for any $\bm{x}$ and $\bm{y}$ (by permuting the region we where we put the base point $z$). 

We are left to consider the case where the algebra element is of length $4$. We claim that for a generator $\bm{x}$ such that $\iota_1\cdot \bm{x}=\bm{x}$, we have $$\langle\tilde{\partial}^2\bm{x}, \rho_{0123}\bm{y}\rangle=\begin{cases}0,\ if\  \bm{x}\neq\bm{y},\\1,\ if\ \bm{x}=\bm{y}.
\end{cases} $$
Assuming this claim, by permuting the subscripts we also have that $\langle\tilde{\partial}^2\bm{x}, \rho_{2301}\bm{y}\rangle$ is $1$ if $\bm{x}=\bm{y}$ and $0$ otherwise, and an idempotent consideration shows $\langle\tilde{\partial}^2\bm{x}, a\bm{y}\rangle=0$ when $a\in\{\rho_{1230},\rho_{3012}\}$. These together imply $\tilde{\partial}^2\bm{x}=\bm{U}\cdot \bm{x}$ when $\iota(\bm{x})=\iota_1$. A similar consideration shows this is true for $\bm{x}$ with $\iota(\bm{x})=\iota_0$ as well. This finishes the proof of the proposition modulo the claim. 

Next, we prove the claim. Note
\begin{equation}\label{Equation, coefficient of extended partial^2}
\langle\tilde{\partial}^2\bm{x}, \rho_{0123}\bm{y}\rangle=\sum_{\bm{w}}\sum_{\substack{\text{ind}(B_i,\overrightarrow{\sigma_i})=1,\\ i=1,2}} \#\mathcal{M}^{B_1}(\bm{x},\bm{w};\overrightarrow{\sigma_1})\#\mathcal{M}^{B_2}(\bm{w},\bm{y};\overrightarrow{\sigma_2}),
\end{equation}
where $(B_i,\overrightarrow{\sigma_i})$ ($i=1,2$) is compatible and $a(-\overrightarrow{\sigma_1})a(-\overrightarrow{\sigma_1})=\rho_{0123}$ or $\bm{U}$; the possible pairs of $(\overrightarrow{\sigma_1},\overrightarrow{\sigma_2})$ are listed below:
\begin{align*}
\bigg\{(\emptyset,\{-\rho_0,-\rho_1,-\rho_2,-\rho_3\}),(\{-\rho_0\},\{-\rho_1,-\rho_2,-\rho_3\}),(\{-\rho_0,-\rho_1,-\rho_2\},\{-\rho_3\}),\\
(\{-\rho_0,-\rho_1\},\{-\rho_2,-\rho_3\}),(\{-\rho_0,-\rho_1,-\rho_2,-\rho_3\},\emptyset),(\emptyset,\{-\rho_0,-\rho_{123}\}),\\(\{-\rho_0\},\{-\rho_{123}\}),(\{-\rho_0,-\rho_{123}\},\emptyset),(\emptyset,\{-\rho_{012},-\rho_{3}\}),(\{-\rho_{012}\},\{-\rho_{3}\}),\\(\{-\rho_{012},-\rho_{3}\},\emptyset),(\emptyset,\{U\}),(\{U\},\emptyset)\bigg\}.
\end{align*}
Let $$\overline{\mathcal{M}}_0\coloneqq\cup_{\text{ind}(B,\overrightarrow{\sigma})=2}\overline{\mathcal{M}}^B(\bm{x},\bm{y};\overrightarrow{\sigma}),$$ where $\overrightarrow{\sigma}\in \{(-\rho_0,-\rho_1,-\rho_2,-\rho_3),(-\rho_0,-\rho_{123}),(-\rho_{012},-\rho_{3})\}$. Let $$\overline{\mathcal{M}}_1:=\cup_{\text{ind}(B,U)=2}\overline{\mathcal{M}}^B(\bm{x},\bm{y};U).$$

Equation \ref{Equation, coefficient of extended partial^2} and the gluing result imply that $\langle\tilde{\partial}^2\bm{x}, \rho_{0123}\bm{y}\rangle$ is equal to the number of two-story ends of the moduli space $\overline{\mathcal{M}}_0\cup \overline{\mathcal{M}}_1$. 

According to Proposition \ref{Proposition, ends of 1-dimensional moduli spaces of 0-P curves, Reeb element of length 4}, the other elements in $\partial \overline{\mathcal{M}}_0$ are: 
\begin{itemize}
\item[(A-1)]simple holomorphic combs with a single split component;
\item[(A-2)]simple boundary degenerations with one corner;
\item[(A-3)]simple boundary degenerations without corners.
\end{itemize}
Proposition \ref{Proposition, ends of moduli space of 1-P curves} shows the other boundary points in $\overline{\mathcal{M}}_1$ in addition to two-story ends are: 
\begin{itemize}
\item[(B-1)]simple holomorphic combs with an orbit curve;
\item[(B-2)]simple boundary degenerations with one corner;
\item[(B-3)]simple boundary degenerations without corners.
\end{itemize}
Note by Proposition \ref{Proposition, ends of 1-dimensional moduli spaces of 0-P curves, Reeb element of length 4} and Proposition \ref{Proposition, ends of moduli space of 1-P curves}, the number of boundary points of type (A-1) is equal to that of type (B-1), both of which is $$\sum_{\text{ind}(B,-\rho_{3012})=1} \#\mathcal{M}^B(\bm{x},\bm{y};-\rho_{3012})+\sum_{\text{ind}(B,-\rho_{1230})=1} \#\mathcal{M}^B(\bm{x},\bm{y};-\rho_{1230}).$$
The parity of the number of boundary points of type (A-2) is equal to that of type (B-2), which are both mod $2$ equal to $$\sum_ {q}\sum_{\{(B_1,B_2)|B_2\in T(q),\ \text{ind}(B_1+B_2;U)=2\}} \#\mathcal{M}^{B_1}(\bm{x},\bm{y};q),$$
where $q$ ranges over self-intersection points of $\alpha_{im}$, and $T(q)$ denotes the set of stabilized teardrops at $q$. 

The parity of the number of boundary points of type (B-3) is even according to Proposition \ref{Proposition, ends of moduli space of 1-P curves}. 

In summary, the parity of the number of boundary points of $\overline{\mathcal{M}}_0\cup \overline{\mathcal{M}}_1$ corresponding to two-story ends is equal to that of type (A-3), which is odd if and only if $\bm{x}=\bm{y}$ by Proposition \ref{Proposition, ends of 1-dimensional moduli spaces of 0-P curves, Reeb element of length 4}. Therefore, $\langle\tilde{\partial}^2\bm{x}, \rho_{0123}\bm{y}\rangle$ is odd if and only if $\bm{x}=\bm{y}$, finishing the proof of the claim. 

In the presence of non-trivial local systems, we simply need to consider the above argument for each domain. For a domain $B$, let $\overline{\mathcal{M}}_0^B$ be the subset of $\overline{\mathcal{M}}_0$ consisting of holomorphic curves with domain $B$, and similarly define $\overline{\mathcal{M}}_1^B$. The two-story ends in $\overline{\mathcal{M}}_0^B\cup \overline{\mathcal{M}}_1^B$ all correspond to the same parallel transport. When ends of type (A-3) do not occur, the two-story ends cancel in pairs by the same argument as above. When ends of type (A-3) appear, we have $B=[\Sigma]$ and $\sigma=(-\rho_0,-\rho_1,-\rho_2,-\rho_3)$. In particular, $\partial_{\alpha_{im}}B=\emptyset$, which induces the identity endomorphism of $\mathcal{E}|_{\bm{x}}$. Also, the number of two-story ends is odd as the number of (A-3) ends is odd. The claim follows from these. 
\end{proof}

\end{document}
